%% =====================================================================
%%  2 / 7   Literature Review  文献综述                         成员 2
%%  三条主线 + 研究空白
%% =====================================================================
\section{Literature Review}
\label{sec:litreview}

\paragraph{Stream 1: from decision support to agentic systems.}
\citet{topol2019} framed the field as a convergence of human and artificial intelligence, in which deep learning would improve diagnosis without displacing the clinician. That framing held while clinical systems remained classifiers. The survey of \citet{wang2025baymax} marks the change: it organizes medical agents by level of involvement and by paradigm, and makes explicit that the object of study is now a system that acts over time rather than a function that returns a label. Concrete instances followed quickly. Agent Hospital simulates an entire hospital in which patient agents and doctor agents evolve together, so that policies can be improved without risk to real patients \citep{li2024agenthospital}. \citet{zhuang2025learning} search over medical agent architectures and report task success of 91.0\% on their benchmark, rising to 98.0\% with a memory module. \citet{jiang2025medagentbench} build a virtual health record in which agents must read, write and act. \citet{schmidgall2024agentclinic} place agents in simulated consultations, \citet{gao2026medexagent} extend evaluation to clinical examination, and \citet{ma2026activemedagent} show that on 175 cases an active agent reaches the correct diagnosis while requesting fewer examinations than a passive one.

The limitation is uniform across the stream: every one of these studies measures what the system can do. None measures who is permitted to let it.

\paragraph{Stream 2: autonomy, oversight and accountability.}
A second literature treats autonomy as a graded property. \citet{ye2025autonomy} propose four levels, from no autonomy through recommendation and action-with-confirmation to full autonomy, and evaluate models at each. This is a genuine advance, because it makes the object of evaluation explicit. But a level describes capability rather than permission: it says what a system is able to do unsupervised, not what it should be allowed to do, nor who answers when it does it. \citet{habli2020accountability} confront the consequence directly, arguing that accountability in AI-supported care cannot be reduced to the performance of a model and must be designed into the sociotechnical system around it. \citet{clayton2025clinicians} defends the clinician's diagnostic role as something that cannot be outsourced without loss, and \citet{abdulnour2025education} turns the point into a curriculum, proposing that supervision of AI use be taught as a clinical skill in its own right.

The gap here is structural. The levels literature lacks an account of authorization; the accountability literature lacks an operational form. Neither tells a clinician what to write down before allowing an agent to act.

\paragraph{Stream 3: collaboration and its costs.}
A third stream studies the human side, empirically and conceptually. \citet{ong2026flight} import aviation's flight rules into clinical AI and argue that collaboration should be governed by explicit, rehearsed procedures rather than by goodwill. \citet{busch2025attitudes} survey hospital patients across several countries and find that acceptance depends on the task, the degree of human involvement and the clinical stakes, rather than on AI in the abstract. The costs are documented too. In a multicentre observational study, endoscopists exposed to AI-assisted colonoscopy showed an absolute decline of 6.0 percentage points in adenoma detection once the assistance was withdrawn \citep{budzyn2025deskilling}; \citet{natali2025deskilling} generalize the finding into a taxonomy of AI-induced deskilling; and \citet{hunter2024reliance} propose reliance drills as a way of measuring dependence before it becomes dangerous. The lineage runs back to \citet{bainbridge1983}, whose ironies of automation --- that automating a task makes the human's residual role harder and less practised --- describe the present situation with uncomfortable precision. \citet{mollick2024} reframes the relationship as co-intelligence, in which the human sets direction and the machine supplies capacity.

What this stream supplies is a vocabulary of failure. What it does not supply is a rule: it can say when collaboration goes wrong, but not what a legitimate delegation consists of.

\paragraph{Research gap.}
Taken together, the three streams leave a specific hole. One measures capability, one grades autonomy, one catalogues the costs of collaboration. None specifies what a delegation must contain. If a clinician is to authorize an agent to act, four things are at stake and none is currently a design requirement. The scope of the authorization must be stated in advance. The authorization must be revocable while the action is in flight. The action must be reconstructable afterwards. And the human's unaided ability must survive the arrangement. The literature can describe each of these failures separately. It has no account of them as a single object --- and that is the object this paper constructs.
